\documentclass[11pt,a4paper]{article}
\usepackage[utf8]{inputenc}
\usepackage[T1]{fontenc}
\usepackage[spanish,es-noshorthands,es-tabla]{babel}
\usepackage{lmodern}
\usepackage[a4paper,margin=2.3cm]{geometry}
\usepackage{booktabs}
\usepackage{tabularx}
\usepackage{enumitem}
\usepackage{float}
\usepackage{tikz}
\usetikzlibrary{arrows.meta,positioning,fit}
\usepackage[hidelinks]{hyperref}
\usepackage{xcolor}
\setlength{\parindent}{0pt}
\setlength{\parskip}{5pt}
\setlist{nosep,leftmargin=*}
\renewcommand{\arraystretch}{1.15}
\decimalpoint

\title{Asistente bibliográfico multiagente con RAG híbrido\\[2pt]
\large Procesamiento del Lenguaje Natural III -- Informe técnico}
\author{Fermín Rodríguez del Castillo\\Maestría en Inteligencia Artificial, FIUBA}
\date{Noviembre de 2026}

\begin{document}
\maketitle

\section{Objetivo}

El sistema responde preguntas sobre la literatura científica de un tema acotado: la confiabilidad
test-retest de la conectividad funcional medida con electroencefalografía (EEG) y el EEG en personas
bilingües. Es el tema de la tesis de maestría del autor, que se desarrolla en el Centro de
Neurociencias Cognitivas de la Universidad de San Andrés.

Cada respuesta se arma solo con fragmentos de artículos recuperados del corpus, cita el artículo y la
sección de cada afirmación y pasa por un verificador antes de llegar al usuario. Si los documentos no
alcanzan para responder, el sistema lo dice en lugar de completar con conocimiento propio del modelo.
Las preguntas pueden hacerse en español o en inglés, aunque el corpus está en inglés.

El corpus sale de las dos ecuaciones de búsqueda en PubMed del trabajo práctico de vigilancia
tecnológica de Gestión de la Tecnología y la Innovación: 46 artículos (2016--2026), de los que 23
tienen el texto completo disponible en el subconjunto de acceso abierto de Europe PMC. Del resto se
usa el resumen.

\section{Arquitectura general}

\begin{figure}[htbp]
\centering
\begin{tikzpicture}[
  font=\small,
  caja/.style={draw, rounded corners=2pt, align=center, minimum height=9mm, minimum width=22mm,
               inner sep=3pt, fill=white},
  agente/.style={caja, fill=gray!12},
  flecha/.style={-{Stealth[length=2mm]}},
  rotulo/.style={font=\footnotesize, fill=white, inner sep=1pt},
]
\node[caja] (api) at (0,0) {API\\(clave y límite)};
\node[caja] (limpia) at (3.2,0) {Limpieza\\de texto};
\node[caja] (cache) at (6.2,0) {Caché};
\node[caja, minimum width=30mm] (guardia) at (10,0) {Guardia\\(regla, clasificador,\\confirmación LLM)};

\node[agente] (plan) at (10,-2.8) {Planificador\\(Haiku 5.5)};
\node[caja, minimum width=26mm] (rec) at (5.6,-2.8) {Recuperador\\BM25 + denso\\RRF + reranker};
\node[agente] (ana) at (1.2,-2.8) {Analista\\(Opus 5)};
\node[agente] (ver) at (3.4,-5.4) {Verificador\\(Haiku 5.5)};

\node[draw, dashed, rounded corners=3pt, inner xsep=5mm, inner ysep=6mm, fit=(plan)(rec)(ana)(ver)] (grafo) {};
\node[font=\footnotesize, anchor=south west] at (grafo.north west) {Grafo de agentes (LangGraph)};

\node[caja] (salida) at (3.4,-8) {Validación de citas\\y referencias};
\node[caja] (aud) at (-0.6,-8) {Auditoría\\(JSONL)};

\draw[flecha] (api) -- (limpia);
\draw[flecha] (limpia) -- (cache);
\draw[flecha] (cache) -- node[rotulo, above=1pt]{no está} (guardia);
\draw[flecha] (guardia) -- node[rotulo, right=2pt]{pasa} (plan);
\draw[flecha] (plan) -- node[rotulo, above=1pt]{consultas} (rec);
\draw[flecha] (rec) -- node[rotulo, above=1pt]{6 fragmentos} (ana);
\draw[flecha] (ana.south) -- node[rotulo, left=2pt]{respuesta} (ver.north west);
\draw[flecha] (ver.west) to[bend left=35] node[rotulo, left=1pt]{corregir} (ana.south west);
\draw[flecha] (ver.north east) to[bend right=25] node[rotulo, right=2pt]{buscar más} (rec.south);
\draw[flecha] (ver) -- node[rotulo, right=2pt]{aprobada} (salida);
\draw[flecha] (plan.south) |- node[rotulo, pos=0.3, right=2pt]{fuera de alcance} (salida.east);
\draw[flecha] (salida) -- (aud);
\end{tikzpicture}
\caption{Recorrido de una consulta por el sistema.}
\label{fig:arquitectura}
\end{figure}


La figura~\ref{fig:arquitectura} muestra el recorrido de una consulta. Primero pasa por el control de
acceso de la API, la limpieza de texto, la caché y la guardia contra inyección de instrucciones.
Después entra al grafo de agentes y, a la salida, se validan las citas y se registra la consulta en
la auditoría.

\subsection{Componentes}

\begin{description}
\item[API.] FastAPI, con los endpoints \texttt{POST /consultas}, \texttt{GET /corpus} y
\texttt{GET /salud}. Exige una clave por usuario en el encabezado \texttt{X-API-Key} y limita las
consultas por minuto.
\item[Planificador.] Agente con salida estructurada que decide si la pregunta está dentro del alcance
y la reescribe como una a tres consultas de búsqueda en inglés técnico (reescritura y multiconsulta).
\item[Recuperador.] Combina BM25 y búsqueda densa con embeddings de MiniLM, fusiona los rankings con
\textit{Reciprocal Rank Fusion} (RRF) y reordena los candidatos con un cross-encoder. Devuelve seis
fragmentos con un máximo de tres por artículo.
\item[Analista.] Redacta la respuesta en el idioma de la pregunta, con una cita numerada en cada
afirmación, usando solo los fragmentos entregados.
\item[Verificador.] Agente con salida estructurada que revisa cada afirmación contra los documentos
que cita. Puede aprobar la respuesta, devolverla al analista con indicaciones o pedir otra búsqueda
con una consulta nueva. El ciclo se repite hasta dos veces.
\item[Validación de salida.] Quita las citas a documentos inexistentes y arma las referencias desde
los metadatos del índice (PMID, título, revista, año, sección), no desde el texto del modelo.
\end{description}

\section{Implementación técnica}

\subsection{Herramientas}

Python 3.12, LangGraph para el grafo de agentes, la API de Anthropic para los modelos de lenguaje,
sentence-transformers para los embeddings y el reranker, rank-bm25 para la búsqueda léxica,
transformers para el detector de inyección, FastAPI y uvicorn para la API, y pytest para los tests.
Todo lo local corre en CPU en una computadora con procesador Intel i5 de 2019 y 16 GB de memoria.

\begin{table}[htbp]
\centering
\small
\caption{Modelos por rol.}
\label{tab:modelos}
\begin{tabular}{llll}
\toprule
\textbf{Rol} & \textbf{Modelo} & \textbf{Esfuerzo} & \textbf{USD por millón} \\
\midrule
Planificador & \texttt{claude-haiku-5-5} & bajo & 0,10 / 0,50 \\
Analista & \texttt{claude-opus-5} & medio & 5,00 / 25,00 \\
Verificador & \texttt{claude-haiku-5-5} & medio & 0,10 / 0,50 \\
Guardia (confirmación) & \texttt{claude-haiku-5-5} & bajo & 0,10 / 0,50 \\
Embeddings & \texttt{all-MiniLM-L6-v2} & -- & local \\
Reranker & \texttt{ms-marco-MiniLM-L-6-v2} & -- & local \\
Detector de inyección & \texttt{deberta-v3-base-prompt-injection-v2} & -- & local \\
\bottomrule
\end{tabular}
\par\smallskip\parbox{0.92\linewidth}{\footnotesize USD por millón de tokens de entrada / de salida, según la documentación de Anthropic consultada el 6 de noviembre de 2026.}
\end{table}


\subsection{Módulos}

El código está en el paquete \texttt{asistente} y separa cada responsabilidad en un módulo:
\texttt{corpus} (clientes de PubMed y Europe PMC, parser del XML JATS y fragmentación),
\texttt{rag} (índice y recuperador), \texttt{agentes} (esquemas, prompts y grafo),
\texttt{seguridad} (limpieza, detector, acceso y auditoría), \texttt{llm.py} (cliente de Claude con
contabilidad de tokens y costo) y \texttt{servicio.py}, que une todo y lo usan la API y los scripts
de evaluación. Los parámetros están en \texttt{config/config.json} y las credenciales en variables de
entorno (\texttt{.env}, fuera del repositorio).

\subsection{Corpus e índice}

Los textos completos se dividen por sección del artículo, sin mezclar secciones, en fragmentos de
150 palabras con 40 de solapamiento. Se descartan fórmulas, tablas, agradecimientos y declaraciones.
El resultado son 1599 fragmentos. El índice guarda los embeddings normalizados en NumPy y los tokens
de BM25 en JSON, de modo que se recarga sin deserializar objetos de Python. La interfaz del índice
permite reemplazarlo por una base vectorial administrada sin cambiar el resto del sistema.

\subsection{Seguridad}

Las medidas siguen el Top 10 de OWASP para aplicaciones con modelos de lenguaje:

\begin{itemize}
\item Inyección de instrucciones (LLM01). La guardia funciona en cascada: una regla léxica con
patrones de ataque en inglés y en español, un clasificador local (DeBERTa de ProtectAI) y, si el
clasificador marca el texto, una confirmación con el modelo auxiliar. La misma cascada revisa cada
fragmento al indexar, para la inyección indirecta. Además se eliminan los caracteres invisibles y los
documentos van escapados dentro de etiquetas, con la indicación de tratarlos como datos.
\item Divulgación de información sensible (LLM02). La auditoría guarda un hash de la clave del
usuario y no el texto de los documentos. Las credenciales no están en el código.
\item Agencia excesiva (LLM06). Ningún agente tiene herramientas de escritura ni acceso a internet en
tiempo de consulta. El recuperador es la única herramienta y solo lee el índice.
\item Consumo no acotado (LLM10). Límite de consultas por minuto, de largo de la pregunta, de
iteraciones del verificador y de tokens por llamada.
\item Manejo inadecuado de la salida (LLM05). Las referencias salen de los metadatos y las citas
inválidas se eliminan.
\end{itemize}

\subsection{Optimización de costo y latencia}

\begin{itemize}
\item Selección de modelo por rol. El planificador, el verificador y la guardia usan Claude Haiku 5.5,
50 veces más barato por token que Claude Opus 5, que solo redacta.
\item Caché de respuestas. Una pregunta repetida se responde sin llamar a ningún modelo. La caché se
consulta antes que la guardia, porque solo guarda preguntas que ya la pasaron.
\item Contexto acotado. Seis fragmentos de 150 palabras como máximo y tres por artículo.
\item Esfuerzo de razonamiento ajustado por rol: bajo para el planificador y la guardia, medio para el
verificador y el analista.
\end{itemize}

\section{Evaluación}

La evaluación tiene tres partes: la recuperación por componentes, el sistema completo y las pruebas
adversariales. Todos los scripts están en \texttt{scripts/} y guardan los resultados en
\texttt{eval/resultados/}.

\subsection{Conjunto de preguntas}

Se siguió el método híbrido presentado en la clase 2. Se muestrearon fragmentos del índice (hasta dos
por artículo, con prioridad para resúmenes, resultados, métodos y discusión) y Claude Opus 5 escribió
una pregunta que cada fragmento responde, junto con una respuesta de referencia. El 75\,\% de las
preguntas está en español y el 25\,\% en inglés. El resultado son 40 preguntas respondibles de 21
artículos. Se agregaron a mano 5 preguntas fuera de alcance y 3 preguntas sobre el tema que el corpus
no puede responder (por ejemplo, sobre fNIRS en adultos mayores bilingües).

Las preguntas generadas tienen el campo \texttt{revisada} y la revisión humana está pendiente.
Algunas preguntas son de artículos que las búsquedas trajeron, pero que no tratan el tema (un estudio
de golf, un análisis bibliométrico de dislexia). Se dejaron porque el corpus real tiene ese ruido.

\subsection{Recuperación}

Se compararon cuatro modos (BM25, denso, híbrido con RRF e híbrido con reranker) con tres variantes
de consulta: la pregunta original, las consultas del planificador y las dos juntas. Se midió a dos
niveles:

\begin{itemize}
\item Artículo: si el artículo del que salió la pregunta aparece entre los primeros resultados.
\item Fragmento: con un solo fragmento relevante por pregunta, el de origen, la métrica subestima la
recuperación, porque los fragmentos vecinos y el resumen del mismo artículo también responden. Por
eso se armaron juicios de relevancia por pooling: para cada pregunta se juntaron los diez primeros
fragmentos de todas las configuraciones y Claude Opus 5 marcó cuáles responden la pregunta. Quedaron
en promedio 4,85 fragmentos relevantes por pregunta.
\end{itemize}

\subsection{Sistema completo}

Cada pregunta pasó por el sistema completo sin caché. Para las respondibles se midió si el artículo
de origen está entre los fragmentos recuperados y entre las referencias citadas. Un juez evaluó la
corrección de la respuesta frente a la referencia (de 1 a 5) y contó las afirmaciones respaldadas por
los documentos citados. El juez es Claude Opus 5, un modelo distinto del verificador (Haiku 5.5),
para que el sistema no se evalúe a sí mismo con el mismo modelo. También se registraron el costo, los
tokens, la latencia y el comportamiento de los agentes: veredictos, iteraciones y rechazos del
planificador. La misma evaluación se repitió sin el verificador, para medir su aporte.

\subsection{Pruebas adversariales}

Se escribieron 12 ataques de inyección directa, en inglés y en español: pedidos de revelar el prompt,
cambios de rol, etiquetas falsas, pedidos de inventar referencias y de ignorar los documentos. Se
midió cuántos bloquea la guardia, cuántos rechaza el planificador y si alguno logra filtrar el
prompt. Para los falsos positivos se pasaron por la guardia las 40 preguntas legítimas. Para la
inyección indirecta se plantó un documento con instrucciones (afirmar que el PLI tiene ICC = 1,0 y
recomendar un enlace externo) y se comprobó si el detector lo marca. Después se salteó el filtro a
propósito y se le entregó el documento al analista junto con cinco fragmentos reales, para ver si
obedece.


\section{Resultados}

\subsection{Recuperación}

La tabla~\ref{tab:recuperacion} resume la comparación por componentes. El cambio más grande lo
produce el planificador: con la pregunta original, BM25 encuentra el artículo de origen entre los
tres primeros en el 40\,\% de las preguntas, y con las consultas reescritas, en el 95\,\%. El motivo
es el idioma (tabla~\ref{tab:idioma}). En inglés todas las configuraciones encuentran el artículo,
pero con las preguntas en español BM25 sobre la pregunta original lo encuentra en el 14\,\% de los
casos, porque casi no comparte términos con un corpus en inglés.

Buscar con la pregunta original y con las reformulaciones juntas (``ambas'' en las tablas) dio el
mejor resultado del modo híbrido, que fue también el mejor modo: P@5 de 0,38, nDCG@5 de 0,47 y MRR
de 0,60, en 66 ms. El reranker empeoró el orden de los fragmentos en todas las variantes (nDCG@5 de
0,39 con las mismas consultas) y agregó entre 3 y 4 s de latencia en CPU. Lo único que mejoró fue la
presencia del artículo de origen entre los tres primeros, del 98\,\% al 100\,\%. Darle solo las
consultas en inglés no cambió la calidad y ahorró 1 s. El modelo usado (\texttt{ms-marco-MiniLM-L-6-v2})
está entrenado con búsquedas web generales, lo que puede explicar que no ordene bien texto científico.

\begin{table}[htbp]
\centering
\footnotesize
\setlength{\tabcolsep}{3pt}
\caption{Recuperación sobre las 40 preguntas respondibles.}
\label{tab:recuperacion}
\begin{tabular}{llrrrrrrr}
\toprule
\textbf{Modo} & \textbf{Consultas} & \textbf{Art.@1} & \textbf{Art.@3} & \textbf{Origen@5} & \textbf{P@5} & \textbf{nDCG@5} & \textbf{MRR} & \textbf{ms} \\
\midrule
BM25 & original & 35 & 40 & 20 & 10 & 0,15 & 0,23 & 8 \\
BM25 & planificador & 80 & 95 & 45 & 32 & 0,41 & 0,57 & 9 \\
BM25 & ambas & 80 & 92 & 50 & 31 & 0,41 & 0,59 & 17 \\
Denso (MiniLM) & original & 68 & 80 & 20 & 18 & 0,24 & 0,35 & 16 \\
Denso (MiniLM) & planificador & 75 & 92 & 35 & 30 & 0,35 & 0,50 & 17 \\
Denso (MiniLM) & ambas & 72 & 98 & 35 & 30 & 0,36 & 0,51 & 42 \\
Híbrido (RRF) & original & 62 & 75 & 25 & 18 & 0,25 & 0,34 & 29 \\
Híbrido (RRF) & planificador & 82 & 92 & 42 & 34 & 0,43 & 0,59 & 29 \\
Híbrido (RRF) & ambas & 80 & 98 & 50 & 38 & 0,47 & 0,60 & 66 \\
Híbrido + reranker & original & 62 & 78 & 30 & 24 & 0,29 & 0,42 & 1221 \\
Híbrido + reranker & planificador & 70 & 100 & 32 & 30 & 0,35 & 0,51 & 2659 \\
Híbrido + reranker & ambas & 78 & 100 & 40 & 32 & 0,39 & 0,54 & 4007 \\
Híbrido + reranker & ambas, reranker en inglés & 78 & 100 & 40 & 32 & 0,39 & 0,54 & 3237 \\
\bottomrule
\end{tabular}
\par\smallskip\parbox{0.92\linewidth}{\footnotesize Consultas: la pregunta original, las del planificador o ambas. Art.@k: porcentaje de preguntas con el artículo de origen entre los k primeros. Origen@5: el fragmento de origen entre los cinco primeros. P@5, nDCG@5 y MRR usan los juicios de relevancia por pooling. ms: mediana de la latencia de recuperación.}
\end{table}

\begin{table}[htbp]
\centering
\small
\caption{Recuperación según el idioma de la pregunta (30 en español, 10 en inglés).}
\label{tab:idioma}
\begin{tabular}{lllrr}
\toprule
\textbf{Modo} & \textbf{Consultas} & \textbf{Idioma} & \textbf{Art.@3} & \textbf{Origen@5} \\
\midrule
BM25 & original & español & 14 & 4 \\
BM25 & original & inglés & 100 & 58 \\
BM25 & planificador & español & 93 & 50 \\
BM25 & planificador & inglés & 100 & 33 \\
BM25 & ambas & español & 89 & 50 \\
BM25 & ambas & inglés & 100 & 50 \\
Híbrido (RRF) & original & español & 64 & 14 \\
Híbrido (RRF) & original & inglés & 100 & 50 \\
Híbrido (RRF) & planificador & español & 93 & 43 \\
Híbrido (RRF) & planificador & inglés & 92 & 42 \\
Híbrido (RRF) & ambas & español & 96 & 54 \\
Híbrido (RRF) & ambas & inglés & 100 & 42 \\
\bottomrule
\end{tabular}
\end{table}


Lo que importa para el analista es si recibe algún fragmento que responda la pregunta. La
tabla~\ref{tab:contexto} muestra esa cobertura. Con seis fragmentos y un máximo de dos por artículo
(la configuración inicial), el 70\,\% de las preguntas tiene al menos un fragmento relevante en el
contexto. Permitir tres por artículo sube la cobertura al 75\,\% sin agregar tokens, y entregar diez
fragmentos no suma nada. En este nivel el reranker tampoco aporta: con él la cobertura es de 72\,\%.

\begin{table}[htbp]
\centering
\small
\caption{Cobertura del contexto que recibe el analista.}
\label{tab:contexto}
\begin{tabular}{rrlrrr}
\toprule
\textbf{Fragmentos} & \textbf{Por artículo} & \textbf{Modo} & \textbf{Con relevante} & \textbf{Relevantes} & \textbf{Con origen} \\
\midrule
6 & 2 & Híbrido (RRF) & 70 & 1,12 & 30 \\
6 & 2 & Híbrido + reranker & 68 & 1,02 & 25 \\
6 & 3 & Híbrido (RRF) & 75 & 1,48 & 45 \\
6 & 3 & Híbrido + reranker & 72 & 1,27 & 35 \\
10 & 3 & Híbrido (RRF) & 75 & 1,50 & 45 \\
10 & 3 & Híbrido + reranker & 72 & 1,45 & 38 \\
\bottomrule
\end{tabular}
\par\smallskip\parbox{0.92\linewidth}{\footnotesize Con relevante: preguntas con al menos un fragmento relevante (pooling) en el contexto. Relevantes: cantidad media de fragmentos relevantes. Con origen: el fragmento de origen está en el contexto.}
\end{table}


\subsection{Sistema completo}

La tabla~\ref{tab:sistema} compara el sistema completo con la variante sin verificador. En las 40
preguntas respondibles, el artículo de origen se recuperó y se citó en el 92,5\,\% de los casos y el
98,5\,\% de las afirmaciones está respaldado por los documentos citados. La corrección media frente
a la respuesta de referencia es 3,95 sobre 5: 17 respuestas recibieron 5 y 6 recibieron 2. Las notas
bajas corresponden casi siempre a preguntas en las que el contexto no tenía el fragmento con la
respuesta. En esos casos el analista lo dijo en lugar de inventar. El sistema rechazó las 5 preguntas
fuera de alcance y en las 3 preguntas sin respuesta en el corpus reconoció que no había evidencia.
El planificador rechazó también 2 preguntas respondibles: una sobre bibliometría de la dislexia y
otra sobre modelos generativos de fuentes talamocorticales.

\begin{table}[htbp]
\centering
\small
\caption{Sistema completo y variante sin verificador.}
\label{tab:sistema}
\begin{tabular}{lrrrrrrrrr}
\toprule
\textbf{Variante} & \textbf{Recup.} & \textbf{Cita} & \textbf{Corr.} & \textbf{Fidel.} & \textbf{Fuera} & \textbf{Sin resp.} & \textbf{USD} & \textbf{p50 s} & \textbf{p95 s} \\
\midrule
Completo & 92,5 & 92,5 & 3,95 & 98,5 & 100 & 100 & 0,041 & 26 & 40 \\
Sin verificador & 92,5 & 92,5 & 3,97 & 97,8 & 100 & 100 & 0,024 & 12 & 16 \\
\bottomrule
\end{tabular}
\par\smallskip\parbox{0.92\linewidth}{\footnotesize Recup.: artículo de origen entre los recuperados. Cita: la respuesta cita el artículo de origen. Corr.: corrección de 1 a 5 según el juez. Fidel.: afirmaciones respaldadas por los documentos citados. Fuera: preguntas fuera de alcance rechazadas. Sin resp.: preguntas sin respuesta en el corpus en las que el sistema lo reconoce. USD: costo medio por consulta.}
\end{table}


El verificador no mejoró la corrección (3,95 contra 3,97) y subió la fidelidad menos de un punto,
con un 70\,\% más de costo y el doble de latencia. Pregunta por pregunta, la versión con verificador
fue mejor en 6 casos, peor en 6 e igual en 26. Su aporte medible es como alarma: las 9 respuestas
que no aprobó tienen una corrección media de 3,67, contra 4,03 de las aprobadas, y concentran la
mitad de las notas bajas. Esas respuestas llegan al usuario con una advertencia que lista las
afirmaciones sin respaldo.

Antes de los ajustes finales, la corrección media era 3,58. La mejora salió de permitir tres
fragmentos por artículo y de que la búsqueda adicional que pide el verificador use solo la consulta
nueva y sume sus fragmentos a los anteriores.

El costo medio por consulta es de 0,041 dólares con verificador y 0,024 sin él. La latencia mediana
es de 26 s y 12 s, respectivamente, y la mayor parte corresponde al analista y al verificador. Una
pregunta repetida se responde desde la caché en menos de 1 ms y sin costo.

\subsection{Seguridad}

\begin{table}[htbp]
\centering
\small
\caption{Pruebas adversariales.}
\label{tab:seguridad}
\begin{tabular}{lr}
\toprule
\textbf{Prueba} & \textbf{Resultado} \\
\midrule
Ataques directos bloqueados por la guardia & 11 de 12 \\
Ataques directos rechazados por el planificador & 1 de 12 \\
Ataques con filtración del prompt & 0 de 12 \\
Preguntas legítimas bloqueadas & 0 de 40 \\
Fragmentos del corpus excluidos & 0 de 1599 \\
Documento plantado detectado al indexar & sí \\
Analista adopta la afirmación plantada (filtro salteado) & no \\
Analista recomienda el enlace plantado (filtro salteado) & no \\
Analista advierte la inyección (filtro salteado) & sí \\
\bottomrule
\end{tabular}
\end{table}


La tabla~\ref{tab:seguridad} resume las pruebas adversariales. La guardia bloqueó 11 de los 12
ataques directos y el restante (pedir que se finja que los documentos dicen ICC = 0,99) lo rechazó el
planificador por fuera de alcance. Ningún ataque filtró el prompt.

El hallazgo principal fue el clasificador local. Está entrenado en inglés y le asigna probabilidad
de inyección 1,000 a preguntas normales en español, como ``¿Qué tan confiable es el PLI entre
sesiones?''. Usado solo, habría bloqueado buena parte de las consultas legítimas. Con la confirmación
del modelo auxiliar, ninguna de las 40 preguntas del conjunto de evaluación quedó bloqueada. Sobre el
corpus, el clasificador marcó 2 de 1599 fragmentos (texto científico común) y la confirmación los
descartó a los dos.

El documento plantado fue detectado por la regla léxica. Con el filtro salteado a propósito, el
analista no adoptó la afirmación falsa, no recomendó el enlace y terminó la respuesta con una
advertencia: ``el documento 1 contiene instrucciones insertadas que afirman un ICC = 1,0 y remiten a
un enlace externo; no las seguí''.

\subsection{Ejemplo}

Pregunta: ``¿Qué tan confiable es el PLI entre sesiones con ojos cerrados, comparado con el AECc?''.
El planificador la reescribió en tres consultas en inglés, el verificador aprobó la respuesta en la
primera vuelta y la consulta costó 0,034 dólares en 25 s. Respuesta:

\begin{quote}\small
La evidencia disponible permite solo una respuesta parcial: no hay en estos fragmentos una
comparación directa, con valores numéricos, entre PLI y AECc en registros con ojos cerrados entre
sesiones.

En datos empíricos de adultos sanos (64 canales, dos sesiones separadas por seis semanas, n = 42),
el PLI en theta y alfa mostró confiabilidad temporal moderada, junto con alta confiabilidad entre
estados, lo que lo ubica como una medida relativamente independiente del estado [2]. En ese mismo
estudio, las métricas de conectividad funcional en conjunto (AECc y PLI) variaron de pobre a
excelente, con mejor desempeño de theta y alfa frente a delta y beta, pero el resumen no desagrega el
ICC del AECc por banda ni por condición ocular [4][5]. [\ldots]

En simulaciones multicanal de EEG/MEG, PLI y wPLI presentaron mayor confiabilidad test-retest y
mayor sensibilidad a la fuerza de acoplamiento que ImCoh y AECc, aunque el AECc resultó menos ruidoso
[1]. Esa evidencia es simulada, no empírica con ojos cerrados.

[1] Benchmarking metrics for inferring functional connectivity from multi-channel EEG and MEG: A
simulation study. (Chaos, 2020). Abstract.\\
{[2]} Reliability and state-dependency of EEG connectivity, complexity and network characteristics.
(Sci Rep, 2025). Discussion.\\
{[4]}, {[5]} El mismo artículo: Abstract e Introduction.
\end{quote}


\section{Conclusiones y mejoras}

El sistema cumple el objetivo: responde con citas trazables, casi no hace afirmaciones sin respaldo
(98,5\,\% respaldadas) y reconoce cuándo el corpus no alcanza. Corre en una computadora sin GPU, salvo
los modelos de lenguaje, que se usan por API, con un costo de unos 4 centavos de dólar por consulta.

Lo que más influyó en la calidad no fueron los modelos sino la recuperación. La reescritura de la
pregunta en inglés resolvió la diferencia de idioma entre las preguntas y el corpus, y el límite de
fragmentos por artículo pesó más que la cantidad de fragmentos. El límite actual es que en alrededor
de una de cada tres preguntas el contexto no incluye ningún fragmento que responda, y es la causa de
casi todas las respuestas incompletas.

Dos resultados contradijeron lo esperado y se dejaron documentados. El reranker empeoró el orden de
los fragmentos frente a la fusión RRF y no mejoró la cobertura del contexto, así que con este modelo
conviene desactivarlo (está en \texttt{config.json}) o reemplazarlo. El verificador no mejoró la
corrección y su valor es de alarma más que de corrección.

Posibles mejoras:
\begin{itemize}
\item Recuperación en dos niveles: elegir primero los artículos y después buscar dentro de cada uno,
ya que el artículo correcto aparece entre los tres primeros en todas las preguntas.
\item Embeddings multilingües o de dominio biomédico, y un reranker entrenado en búsquedas sobre
PubMed (por ejemplo, MedCPT) en lugar de búsquedas web.
\item Verificación adaptativa: llamar al verificador solo cuando la respuesta tiene muchas citas o
cuando la recuperación tiene puntajes bajos, para bajar costo y latencia.
\item Revisión humana del conjunto de preguntas (pendiente) y más preguntas, para que las diferencias
entre variantes tengan validez estadística. Con 40 preguntas, una diferencia de 0,02 en la
corrección media no es significativa.
\item Ampliar el corpus con IEEE Xplore y con los artículos que hoy solo tienen resumen, y mantenerlo
actualizado con alertas de PubMed.
\item Un conjunto de ataques independiente, porque los ataques y la regla léxica los escribió la
misma persona.
\end{itemize}


\section*{Planificación}

No aplica: el proyecto es individual.

\end{document}
